% !TEX root = tarasov-qasp-zh.tex

\chapter[失重状态下摆会怎样？]{失重状态下摆会怎样？\label{ch-12}}

\begin{dialogue}

\tchr 设想在电梯的厢壁上钉一枚钉子，用长为 $l$ 的细绳把一个摆锤（bob）悬挂在钉子上。然后使摆锤运动起来，让它做简谐振动（harmonic vibrations）。设电梯以向上的加速度 $a$ 上升。摆的振动周期是多少？

\stda 乘加速上升的电梯时，我们会感到重量有所增加。想必摆也会“感觉到”同样的增加。我想，它的振动周期可以由公式
\begin{equation}%
T = 2 \pi \sqrt{\frac{l}{g+a}}
\label{eq-77}
\end{equation}
求出。不过，我还不能足够严格地论证这个公式。

\tchr 你的公式是正确的。但要论证它，我们得采取一种对大家来说全新的观点。到目前为止，我们讨论的只是位于惯性参考系中的物体，一直回避非惯性系；我甚至专门提醒过你们不要使用非惯性参考系（\chap{ch-04}）。尽管如此，在现在这一节里，用这种参考系反而更方便——就本题而言，就是固连于加速电梯的参考系。

回想一下：在以加速度 $a$ 运动的非惯性参考系中考察质量为 $m$ 的物体的运动时，我们必须纯粹形式地给物体再施加一个附加力。这个力叫做惯性力，大小等于 $ma$，方向与加速度的方向相反。给物体加上惯性力之后，我们就可以不再考虑参考系正在加速运动这件事，而把运动当作在惯性系中那样来处理。在电梯这道题里，我们必须给摆锤附加一个力 $ma$。这个力大小恒定、方向不变，并且与重力 $mg$ 的方向一致。由此可知，在式~\eqref{eq-75} 中应把加速度 $g$ 换成两个加速度的算术和 $(g+a)$。结果就得到你提出的公式~\eqref{eq-77}。
\pnote{学生说“感到重量增加”，说的就是人教版必修第一册第四章第 6 节“超重和失重”中的超重：电梯加速上升时，体重秤示数为 $m(g+a)$。式~\eqref{eq-77} 中把 $g$ 换成 $g+a$，来历相同——在电梯参考系里，惯性力与重力同向叠加，摆锤的“视重”变为 $m(g+a)$，回复力相应增大，周期变短。}

\stda 这么说，如果电梯以向下的加速度 $a$ 下降，振动周期就由两个加速度之差 $(g-a)$ 决定了，因为这时惯性力 $ma$ 与重力方向相反。对吗？
\pnote{底本本章有若干拼写级排印小误：上句的惯性力原文误作 ``rna''（应为 $ma$），三处图题的 ``Anaysing''（应为 Analysing），式~\eqref{eq-79} 第二行一处正弦函数名未排正体；中文版均径改，不另注。另：本章式~\eqref{80} 的标签原文即缺 ``eq-'' 前缀（全书无引用处），照底本保留。}

\tchr 当然。这时摆的振动周期为
\begin{equation}%
T = 2 \pi \sqrt{\frac{l}{g-a}}
\label{eq-78}
\end{equation}
这个公式只在 $a<g$ 的条件下才有意义。加速度 $a$ 的值越接近 $g$，摆的振动周期越大。当 $a=g$ 时，便进入失重状态。这时摆会怎样呢？
\pnote{原书“只在 $a<g$ 时有意义”是就摆锤悬在钉子下方的通常位形而言的。若电梯向下的加速度超过 $g$（即 $a>g$），惯性力的大小超过重力：等效重力加速度大小变为 $a-g$、方向竖直向上，摆将改以绳竖直向上、摆锤位于悬点上方的平衡位置振动，周期为 $2\pi\sqrt{l/(a-g)}$。}

\stda 按公式~\eqref{eq-78}，周期会变得无穷大。这必定意味着摆是静止的。

\tchr 我们来把你回答中的一些细节弄清楚。起初，摆在电梯里振动着。突然，电梯脱开悬挂，开始向下自由坠落（忽略空气阻力）。这时摆会怎样？

\stda 就像我刚才说的，摆会停下来。

\tchr 你的回答不完全正确。如果电梯脱开悬挂的那个时刻，摆锤恰好处于它的两个极端位置之一，摆确实会静止（当然是相对于电梯而言）。但如果在那个时刻摆锤不在极端位置上，它就将继续在竖直平面内绕绳端转动，速率保持不变，等于发生这起“事故”的时刻它所具有的速率。
\pnote{失重开始的那一瞬，摆锤的速度本就沿圆弧切向、与绳垂直；此后在电梯参考系中，它只受绳的张力——张力沿绳指向悬点、处处与速度垂直，只改变速度方向而不做功，速率自然不变，所需的向心力 $mv^{2}/l$ 恰由张力提供。匀速圆周运动正由此而来。}

\stda 我明白了。

\tchr 那么，请画一张图，说明处于失重状态的宇宙飞船里摆（系在绳上的摆锤）的行为。

\stda 在宇宙飞船里，绳端的摆锤要么静止（相对于飞船），要么沿一个圆周转动，圆的半径由绳长确定（当然，前提是飞船的墙壁和天花板不碍事）。

\tchr 你描绘的图景还不完全。设我们身处失重状态的宇宙飞船之中。取来摆锤和细绳，把绳的自由端固定住，使墙壁和天花板都不妨碍摆锤的运动。然后轻轻地放开摆锤。小球保持静止。这里要区分两种情形：
\begin{enumerate*}[label=(\arabic*)]
\item 绳是松弛的；
\item 绳是拉紧的。
\end{enumerate*}
\begin{figure}[!h]%
\centering
\includegraphics[width=0.45\textwidth]{figs/ch-12/fig-051a.pdf}
\caption[宇宙飞船中摆的运动分析。]{分析宇宙飞船中摆的运动。\label{fig-51}}
\end{figure}
先看第一种情形（\figr{fig-51}~\drkgry{(a)} 中位置 1）。给摆锤一个速度 $v_{0}$。于是摆锤将做匀速直线运动，直到绳被拉直为止（\figr{fig-51}~\drkgry{(a)} 中位置 2）。在这一时刻，绳将对摆锤施以反作用力，其方式与球从墙上弹开时墙对球的作用相同。结果，摆锤运动的方向发生突变，随后它重新做匀速直线运动（\figr{fig-51}~\drkgry{(a)} 中位置 3）。在这种奇特的“反射”中，入射角与反射角相等的规则应当成立。
\pnote{原书把张紧的绳当作“弹性镜面”处理：沿绳方向的分速度等大反向地弹回，机械能没有损失，摆锤随后沿直线运动、绳重新松弛，如此往复，轨迹是半径为 $l$ 的圆内一串等长的弦。国内中学与竞赛处理“绳突然张紧”用的是另一个极端模型：不可伸长的绳以冲击式张力把沿绳方向的分速度抵消为零（机械能有损失），摆锤此后只保留垂直于绳的分速度、转为沿圆周运动。两者都是理想化极限；原书的模型自洽——正因弹回后绳重新松弛，摆锤才继续走直线——但与国内惯用模型不同，对照阅读时须留意。}

再看第二种情形：先把绳拉紧，然后轻轻地放开摆锤。与第一种情形一样，摆锤将保持在放开时的位置上静止（\figr{fig-51}~\drkgry{(b)} 中位置 1）。随后沿垂直于绳的方向给摆锤一个速度 $v_{0}$。于是摆锤开始沿圆周匀速转动。转动平面由绳与传给摆锤的速度矢量共同确定。
\begin{figure}[!h]%
\centering
\includegraphics[width=0.45\textwidth]{figs/ch-12/fig-052a.pdf}
\caption[摆在加速参考系中的运动分析。]{分析摆在加速参考系中的运动，该参考系的加速度与地球重力加速度成某一角度。\label{fig-52}}
\end{figure}
我们来考察下面的问题。\hlred{长为 $l$ 的细绳，一端系着摆锤，另一端固定在小车上；小车沿倾角为 $\alpha$ 的斜面无摩擦地滑下（\figr{fig-52}~\drkgry{(a)}）。求处于以某一加速度运动的参考系中的这个摆的振动周期。}不过，与前面电梯的各题不同，这里系统的加速度与地球重力加速度成某一角度。这就提出了一个附加问题：\hlred{摆绳的平衡方向如何？}

\stda 这种题我从前试着分析过一次，可是越想越乱，没能解出来。

\tchr 这种情形下摆的振动周期仍由公式~\eqref{eq-75} 确定，只是与电梯的情形一样，须把其中的 $g$ 换成某个等效加速度（effective acceleration）。这个加速度（我们记作 $g_{e\!f\!\!\!f}$）等于重力加速度与给定系统的加速度的矢量和。还要注意，在上述和式中，小车的加速度矢量应以反号出现，因为惯性力沿着与系统加速度相反的方向。各加速度矢量如\figr{fig-52}~\drkgry{(b)} 所示，小车的加速度等于 $g \sin \alpha$。下面求 $g_{e\!f\!\!\!f}$：
\begin{equation}%
\begin{split}
g_{e\!f\!\!\!f} & = \sqrt{g_{e\!f\!\!\!fx}^{2} + g_{e\!f\!\!\!fy}^{2}} \\
& = \sqrt{ (g \sin \alpha \cos \alpha)^{2} + (g-g \sin^{2}\alpha)^{2} }\\
& = g \cos \alpha
\end{split}
\label{eq-79}
\end{equation}
由此得
\begin{equation}%
T = 2 \pi \sqrt{\frac{l}{g \cos \alpha}}
\label{80}
\end{equation}
\pnote{这就是竞赛教辅所说的“等效重力加速度”法：参考系以加速度 $a_{0}$ 平动时，等效重力加速度等于重力加速度与 $a_{0}$ 的矢量差，单摆周期为 $T=2\pi\sqrt{l/g_{e\!f\!\!\!f}}$，平衡位置沿 $g_{e\!f\!\!\!f}$ 方向。电梯各题是二者共线的特例；本题参考系的加速度沿斜面向下、大小为 $g\sin\alpha$，故 $g_{e\!f\!\!\!f}=g\cos\alpha$。此法超出中学课程范围，属物理竞赛与大学普通物理内容。}

\stda 怎样确定绳的平衡方向呢？

\tchr 它就是加速度 $g_{e\!f\!\!\!f}$ 的方向。由式~\eqref{eq-79} 容易看出，这个方向与竖直方向成角 $\alpha$。换句话说，在平衡位置上，沿斜面滑下的小车上的摆，其绳将垂直于斜面。

\stdb 这最后一个结果能不能用别的办法得到？

\tchr 可以。直接考察摆锤相对于小车的平衡，就能得出同样的结论。作用在摆锤上的力有：重力 $mg$、绳的张力 $T$ 和惯性力 $ma$（\figr{fig-53}）。
\begin{figure}[!h]%
\centering
\includegraphics[width=0.4\textwidth]{figs/ch-12/fig-053a.pdf}
\caption[摆的运动分析。]{分析摆的运动。\label{fig-53}}
\end{figure}
用 $\beta$ 表示绳与竖直方向之间的夹角。接着把这些力沿竖直方向和水平方向分解，再对每个方向分别写出分力的平衡条件。于是有
\begin{equation}%
\left.
\begin{aligned}
T \cos \beta +ma \sin \alpha & =mg \\
T \sin \beta & = ma \cos \alpha
\end{aligned}
\right\}
\label{eq-81}
\end{equation}
注意到 $a=g \sin \alpha$，把方程组~\eqref{eq-81} 改写为
\begin{equation*}%
\left.
\begin{aligned}
T \cos \beta & = mg (1 - \sin^{2} \alpha) \\
T \sin \beta & = mg \sin \alpha \cos \alpha
\end{aligned}
\right\}
\end{equation*}
两式相除，得
\begin{equation*}%
\cot \beta = \cot \alpha
\end{equation*}
于是，角 $\beta$ 与 $\alpha$ 相等。因此，摆绳的平衡方向垂直于斜面。

\stdb 我一直紧跟您的讲解，最后得出一个结论：当初回答您关于卫星受力的问题时，我举出引力和离心力（见\chap{ch-08}），也并不算全错。

只是，我的回答应当放到固连于卫星的参考系中去理解，而其中的离心力应理解为惯性力。在固连于卫星的非惯性参考系中，我们要处理的已经不是动力学问题，而是静力学问题：这是一个力的平衡问题，其中一个力就是惯性离心力。

\tchr 用这种观点处理卫星问题是允许的。不过，当初在\chap{ch-08} 里提到离心力时，你并没有把它当作惯性力；你不过是想编造出一个什么东西，好让卫星不落到地球上。况且在你所说的这个情形里，根本没必要转换到固连于卫星的参考系：不引入惯性离心力，问题的物理实质反而展示得更清楚。我以前的忠告现在仍然有效：\hlblue{如果没有特殊需要，就不要使用非惯性参考系。}

\end{dialogue}

\cleardoublepage

\thispagestyle{empty}
\vspace*{-1.5cm}
\begin{figure*}[!ht]
\centering
\includegraphics[width=0.75\textwidth]{figs/sec/tarasov-qasp-sec-07.pdf}
\end{figure*}
\begin{fullwidth}
{\textsf{\Large \hlred{静力学的定律就是平衡的定律。请仔细地学习这些定律。不要忘记，它们具有极大的实际意义。一个不懂得静力学基本定律的建筑工人是不可想象的。我们将考察几个说明力的分解法则的例子。随后的讨论涉及物体的平衡条件，这些条件尤其被用来确定重心的位置。
}}}
\end{fullwidth}

% NEWTERMS: harmonic vibrations | 简谐振动 | 人教版选择性必修第一册作“简谐运动”；第 10 章章尾格言已用“简谐运动”，本章为与全书 vibration=振动 一致取“简谐振动”（与第 10、11 章的冲突声明，全书收尾时统一）
% NEWTERMS: bob | 摆锤 | 本章首现括注（bob）；底本第 11 章以 a ball (called a bob) 正式引入，第 11 章译出后以彼处为全书首现
% NEWTERMS: force of inertia | 惯性力 | 即第 4 章 inertial force，同义异写，不重复括注
% NEWTERMS: state of weightlessness | 失重状态 | 失重已由第 7 章格言、第 9 章正文定名，不重复括注
% NEWTERMS: effective acceleration | 等效加速度 | 即竞赛教辅所称“等效重力加速度”，等于重力加速度与参考系加速度的矢量差；中学不作要求
% NEWTERMS: equilibrium direction | 平衡方向 | 摆绳处于平衡时所指方向，沿等效加速度方向
% NEWTERMS: centrifugal force of inertia | 惯性离心力 | 转动参考系中的惯性力；与第 8 章辨析后弃用的“离心力”相区别
% NEWTERMS: reaction of the string | 绳的反作用力 | 指绳的约束力（张力），俄式教材用语直译
% NEWTERMS: statics | 静力学 | 首现于学生乙谈卫星参考系处；下一章主题
% NEWTERMS: truck | 小车 | 沿斜面下滑的载具，习题语境通用译名
